% Shared Round-2 stress benchmark. Templates provide \BenchCallout.
\section{高强度中文正文：行距、行长与段落节奏}
\label{sec:prose}
一个可靠的 PDF 排版体系，不应该要求作者在每一页上手工调整位置，也不应该依赖大量卡片来制造“设计感”。真正稳定的版面来自一组相互约束的基础参数：版心宽度决定一行能够容纳多少文字，字号与行距共同决定阅读密度，段首缩进和段间距决定段落边界是否清楚，标题前后的留白决定章节层级能否在快速翻阅时被识别，而分页规则则必须保证长文、表格、公式与图形能够自然地跨页。中文正文尤其容易暴露不合理的行距：如果基线距离太小，汉字的方形轮廓会在视觉上形成连续的黑色带状区域；如果行距过大，又会让段落失去整体感。第二轮因此将中文正文的默认行距提高到更接近出版物的范围，并要求在多页连续阅读时仍然舒适。

为了让测试具有足够压力，本段继续增加长度。模板不能只在一页短内容上好看，而必须在真实知识型资料中维持稳定节奏。用户可能连续阅读十几页，期间会遇到公式、表格、插图、引文、脚注和多级标题；页面既可能非常密集，也可能因为章节转换而相对疏朗。好的 LaTeX 模板应该让这些变化看起来是同一本书、同一份报告或同一套手册的一部分，而不是每遇到一种内容就突然切换成另一种网页组件。这也是为什么本轮把正文优先设为硬规则：视觉元素可以帮助理解，但不能抢走正文的主导地位。

再加入一段用于观察中文标点与英文字母混排。XeLaTeX、CTeX、KOMA-Script、memoir、TikZ、PGFPlots、tcolorbox 等组件共同工作时，英文缩写与中文之间的间距应自然。数字如 2026、A4、200 DPI、SHA-256 也应保持统一基线。加入较少见汉字“龘、麤、鱻”用于确认 CJK 字体覆盖；真正缺字必须被严格日志门禁显式拒绝。

\BenchCallout{排印基线}{本轮把正文行距作为独立验收项。目标不是统一到某一个数值，而是确保 10.5--11.5pt 中文正文在 A4 版心内具有足够呼吸感；连续三页阅读时不能形成压迫性的黑色文字带。}

\Needspace{7\baselineskip}
\subsection{一个故意很长的标题：用于验证标题换行后仍然保持层级、不会贴近页底，也不会把紧随其后的正文挤成孤行}
这个标题故意超过常见长度。模板应允许标题自然换行，并在标题与正文之间保持稳定距离。若标题落在页尾，needspace 应帮助它与后续正文一起移动到下一页，从而减少孤标题现象。

\section{公式系统：行内、独立公式与多行对齐}
正文中先出现行内公式 $E_k=\frac12 mv^2$，随后给出动量守恒：
\[
 m_1v_1+m_2v_2=m_1v'_1+m_2v'_2.
\]
再用多行环境测试上下留白和等号对齐：
\begin{align}
 p_{\mathrm{before}} &= m_1v_1+m_2v_2,\\
 p_{\mathrm{after}}  &= m_1v'_1+m_2v'_2,\\
 \Delta p            &= p_{\mathrm{after}}-p_{\mathrm{before}}=0.
\end{align}
如果两物体发生完全非弹性碰撞并粘在一起，则共同速度为
\[
 v=\frac{m_1v_1+m_2v_2}{m_1+m_2}.
\]
公式前后不能出现过短的行距，也不能因为大公式造成突兀的空洞。

\section{长表格：跨页、表头重复与窄列中文换行}
下面的长表格故意跨页。每一页需要重复表头，列宽不能溢出版心，中文换行不能产生极端窄行。
\small
\begin{longtable}{@{}p{.13\linewidth}p{.36\linewidth}p{.39\linewidth}@{}}
\caption{Round 2 跨页长表压力测试}\label{tab:stress}\\
\toprule
编号 & 检查对象 & 合格标准 \\
\midrule
\endfirsthead
\toprule 编号 & 检查对象 & 合格标准 \\
\midrule
\endhead
\bottomrule
\endfoot
R01 & 中文正文与分页组合测试第 01 项 & 行距舒适、列宽稳定、表头重复，且不出现裁切、重叠、异常空白或不可读的窄列。\\
R02 & 中文正文与分页组合测试第 02 项 & 行距舒适、列宽稳定、表头重复，且不出现裁切、重叠、异常空白或不可读的窄列。\\
R03 & 中文正文与分页组合测试第 03 项 & 行距舒适、列宽稳定、表头重复，且不出现裁切、重叠、异常空白或不可读的窄列。\\
R04 & 中文正文与分页组合测试第 04 项 & 行距舒适、列宽稳定、表头重复，且不出现裁切、重叠、异常空白或不可读的窄列。\\
R05 & 中文正文与分页组合测试第 05 项 & 行距舒适、列宽稳定、表头重复，且不出现裁切、重叠、异常空白或不可读的窄列。\\
R06 & 中文正文与分页组合测试第 06 项 & 行距舒适、列宽稳定、表头重复，且不出现裁切、重叠、异常空白或不可读的窄列。\\
R07 & 中文正文与分页组合测试第 07 项 & 行距舒适、列宽稳定、表头重复，且不出现裁切、重叠、异常空白或不可读的窄列。\\
R08 & 中文正文与分页组合测试第 08 项 & 行距舒适、列宽稳定、表头重复，且不出现裁切、重叠、异常空白或不可读的窄列。\\
R09 & 中文正文与分页组合测试第 09 项 & 行距舒适、列宽稳定、表头重复，且不出现裁切、重叠、异常空白或不可读的窄列。\\
R10 & 中文正文与分页组合测试第 10 项 & 行距舒适、列宽稳定、表头重复，且不出现裁切、重叠、异常空白或不可读的窄列。\\
R11 & 中文正文与分页组合测试第 11 项 & 行距舒适、列宽稳定、表头重复，且不出现裁切、重叠、异常空白或不可读的窄列。\\
R12 & 中文正文与分页组合测试第 12 项 & 行距舒适、列宽稳定、表头重复，且不出现裁切、重叠、异常空白或不可读的窄列。\\
R13 & 中文正文与分页组合测试第 13 项 & 行距舒适、列宽稳定、表头重复，且不出现裁切、重叠、异常空白或不可读的窄列。\\
R14 & 中文正文与分页组合测试第 14 项 & 行距舒适、列宽稳定、表头重复，且不出现裁切、重叠、异常空白或不可读的窄列。\\
R15 & 中文正文与分页组合测试第 15 项 & 行距舒适、列宽稳定、表头重复，且不出现裁切、重叠、异常空白或不可读的窄列。\\
R16 & 中文正文与分页组合测试第 16 项 & 行距舒适、列宽稳定、表头重复，且不出现裁切、重叠、异常空白或不可读的窄列。\\
R17 & 中文正文与分页组合测试第 17 项 & 行距舒适、列宽稳定、表头重复，且不出现裁切、重叠、异常空白或不可读的窄列。\\
R18 & 中文正文与分页组合测试第 18 项 & 行距舒适、列宽稳定、表头重复，且不出现裁切、重叠、异常空白或不可读的窄列。\\
R19 & 中文正文与分页组合测试第 19 项 & 行距舒适、列宽稳定、表头重复，且不出现裁切、重叠、异常空白或不可读的窄列。\\
R20 & 中文正文与分页组合测试第 20 项 & 行距舒适、列宽稳定、表头重复，且不出现裁切、重叠、异常空白或不可读的窄列。\\
R21 & 中文正文与分页组合测试第 21 项 & 行距舒适、列宽稳定、表头重复，且不出现裁切、重叠、异常空白或不可读的窄列。\\
R22 & 中文正文与分页组合测试第 22 项 & 行距舒适、列宽稳定、表头重复，且不出现裁切、重叠、异常空白或不可读的窄列。\\
R23 & 中文正文与分页组合测试第 23 项 & 行距舒适、列宽稳定、表头重复，且不出现裁切、重叠、异常空白或不可读的窄列。\\
R24 & 中文正文与分页组合测试第 24 项 & 行距舒适、列宽稳定、表头重复，且不出现裁切、重叠、异常空白或不可读的窄列。\\
R25 & 中文正文与分页组合测试第 25 项 & 行距舒适、列宽稳定、表头重复，且不出现裁切、重叠、异常空白或不可读的窄列。\\
R26 & 中文正文与分页组合测试第 26 项 & 行距舒适、列宽稳定、表头重复，且不出现裁切、重叠、异常空白或不可读的窄列。\\
R27 & 中文正文与分页组合测试第 27 项 & 行距舒适、列宽稳定、表头重复，且不出现裁切、重叠、异常空白或不可读的窄列。\\
R28 & 中文正文与分页组合测试第 28 项 & 行距舒适、列宽稳定、表头重复，且不出现裁切、重叠、异常空白或不可读的窄列。\\
\end{longtable}
\normalsize

\clearpage
\section{图形与图像：TikZ、PGFPlots、真实外部页面资源}
\subsection{TikZ 结构图}
\begin{figure}[htbp]
\centering
\begin{tikzpicture}[>=stealth,thick,node distance=18mm]
\node[draw,rounded corners,fill=blue!5,minimum width=31mm,minimum height=10mm] (a) {内容审查};
\node[draw,rounded corners,fill=blue!5,right=of a,minimum width=31mm,minimum height=10mm] (b) {资源构建};
\node[draw,rounded corners,fill=blue!5,right=of b,minimum width=31mm,minimum height=10mm] (c) {最终组合};
\draw[->] (a)--(b);\draw[->] (b)--(c);
\node[below=6mm of b,align=center,text width=90mm]{每一阶段都必须保留可追踪证据；任何上游哈希变化都会使下游验收失效。};
\end{tikzpicture}
\caption{TikZ 流程图压力测试}
\end{figure}

\subsection{PGFPlots 曲线}
\begin{figure}[htbp]
\centering
\begin{tikzpicture}
\begin{axis}[width=.82\linewidth,height=58mm,xlabel={训练轮次},ylabel={稳定性},grid=major,legend pos=south east]
\addplot[very thick,mark=*] coordinates {(1,42)(2,61)(3,74)(4,83)(5,90)};\addlegendentry{结构化流程}
\addplot[dashed,thick,mark=square*] coordinates {(1,45)(2,49)(3,54)(4,57)(5,59)};\addlegendentry{临时手工调整}
\end{axis}
\end{tikzpicture}
\caption{PGFPlots 图形与图注间距测试}
\end{figure}

\subsection{真实外部页面资源与来源页占位}
图~\ref{fig:sourcefixture} 是先独立编译出的外部 PDF 页面资源，用于验证 \texttt{includegraphics} 缩放、图注和分页。
\begin{figure}[htbp]
\centering
\includegraphics[width=.58\linewidth]{source-page-fixture.pdf}
\caption{中性来源页 PDF fixture}\label{fig:sourcefixture}
\end{figure}

\section{引文、列表、脚注与链接}
\begin{quote}
好的模板不是把每一页做成海报，而是让读者在没有意识到排版存在的情况下持续阅读，并在需要时迅速定位标题、表格、公式和图形。
\end{quote}
\begin{enumerate}[label=\arabic*.,leftmargin=2.2em,itemsep=.35em]
  \item 正文应保持连贯，不因列表切换产生过大的垂直跳跃；
  \item 二级条目需要与正文对齐，而不是越来越向右挤压；
  \item 较长条目换行后，续行应与文字起点对齐；
  \item 脚注\footnote{脚注必须足够清晰，同时不能与正文争夺视觉注意力。}应保持较小字号但仍可读。
\end{enumerate}
测试链接：\url{https://ctan.org/}。长 URL 不能冲出版心。

这一节再加入两段正常长度正文，用来避免图像浮动结束后只剩几行文字的偶然稀疏页，同时观察脚注、列表和普通段落连续出现时的节奏是否稳定。模板不能依赖手工分页来填满页面；它应该让 LaTeX 的正常分页算法自然处理内容。

第二段继续检查链接后的断行与段尾位置。如果当前页尚有空间，这些文字应合理利用剩余版心，而不是留下一个只占页面顶部很小区域的伪空白页。这类问题在压力测试初版中出现过，因此将它作为真实失败修复，而不是解释为设计留白。

\clearpage
\section{故意稀疏的一页：验证留白是否像设计而不是错误}
这一页只保留少量正文和一个核心结论，用来测试模板在低密度内容下是否仍然平衡。稀疏页不应被机器简单判定为失败，但必须在视觉验收中确认它不是意外空白。
\vspace{16mm}
\begin{center}
\Large\bfseries 排印系统首先服务连续阅读，\\[1.2ex]
其次才是视觉装饰。
\end{center}
\vspace{18mm}
\BenchCallout{稀疏页观察}{页眉、页脚、标题和正文应形成稳定的垂直重心；不应出现内容全部挤在上半页、下半页完全失衡的情况。}

\clearpage
\section{最后的密集页：连续正文、交叉引用与综合检查}
根据表~\ref{tab:stress}、图~\ref{fig:sourcefixture} 以及第~\ref{sec:prose} 节的正文测试，本轮最终关注的是系统性稳定而不是单页漂亮。一个模板只有同时处理长文、跨页表格、公式、矢量图、外部页面资源、引文、脚注、链接和目录，才有资格进入第三轮真实业务 PDF 测试。

为了制造密集页面，这里再加入两个连续长段。第一段用于观察正文在页底附近如何断行。LaTeX 通常能够自动完成合理分页，但标题、浮动体与长表格组合时仍可能产生令人不舒服的空白。模板应尽量避免人为的强制定位，使用成熟的浮动机制和适度的 Needspace 约束。第二段继续强调视觉差异度：即使把所有强调色转换成灰度，五个模板仍应该能通过版心、章节结构、边注、信息框、页眉页脚、栏目组织等结构特征被区分；如果只能靠换颜色区分，就不应被视为两个独立模板。

第二段用于确认新的行距基线没有在长文末尾退化。正文的基线距离必须在整个文档一致，不能因为进入表格、列表或图形之后发生不可解释的变化。中文段首缩进应保持统一；只有标题、引文、列表、图注等具有明确语义的环境才允许使用不同的缩进与间距。最终验收会同时参考机器预检、全页渲染和人工视觉检查，任何一项都不能单独代表通过。
